\section{初始失效进程情形}

在完成主定理证明之后，原论文还讨论了一个不同的故障情形：进程可能在协议开始前就已经失效，但协议执行过程中不会再有新的进程死亡。这个部分并不是主定理证明的一环，而是用来说明不可能性结果的适用边界。

主定理考虑的是执行过程中可能发生崩溃的完全异步共识。在这种模型下，某个进程长时间没有行动，其他进程无法区分它是已经崩溃，还是只是运行很慢。正是这种不可区分性，使协议无法同时保证一致性和终止性。初始失效进程情形则改变了故障假设：哪些进程已经失效在执行开始时就固定下来，之后存活进程不再发生新的崩溃。

在这一较弱的故障模型下，原论文给出一个正向结果：如果初始时存活的进程占严格多数，并且执行过程中没有进程继续死亡，那么存在一个部分正确的共识协议，使所有非故障进程最终都能作出决定。这个结论说明，FLP 不可能性并不是说只要系统中存在不可用进程，共识就一定无法解决；它针对的是更困难的动态崩溃情形。

原论文中的协议大致分为两个阶段。第一阶段，每个存活进程广播自己的进程编号，并等待收到足够多其他进程的消息。根据“哪个进程收到了哪个进程的消息”，系统可以形成一个有向图。图中的节点对应进程，如果进程 $j$ 收到了进程 $i$ 的消息，就在图中加入从 $i$ 到 $j$ 的边。

第二阶段，进程继续交换自己已经知道的图信息和初始值信息，从而逐步获得关于图的传递可达关系。每个完成协议的进程最终都能识别出同一个初始团。这里的初始团可以理解为一个没有外部入边的共同进程集合。由于初始存活进程占严格多数，论文证明这样的初始团是唯一的，并且每个完成第二阶段的进程都能知道这个集合及其成员的初始值。

最后，所有进程根据这个共同初始团中进程的初始值，使用同一个预先约定的规则作出决定。由于完成协议的进程识别出的初始团相同，且它们使用的决定规则也相同，所以它们会得到相同的决定值。这保证了一致性；同时，在该模型假设下，非故障进程能够完成两个阶段，因此也能最终决定。

这一部分和主定理的关系可以概括为对照关系。主定理证明的是：在完全异步系统中，如果允许进程在执行过程中崩溃，则不存在确定性协议能够保证所有容许执行都终止。初始失效进程部分则说明：如果把故障限制为协议开始前已经发生，并要求执行过程中没有新的崩溃，那么在多数进程存活的条件下，共识仍然可以解决。因此，它帮助界定了 FLP 结果的真正含义：困难不只是“有进程不可用”，而是“进程可能在异步执行过程中停止，且其他进程无法判断它是慢还是死”。
